% Archived theory-oriented draft; the SE-oriented paper is main.tex.
\documentclass[10pt]{article}

\usepackage[a4paper,margin=0.78in]{geometry}
\usepackage{amsmath,amssymb,amsthm}
\usepackage{booktabs}
\usepackage{enumitem}
\usepackage{graphicx}
\usepackage{listings}
\usepackage{xcolor}
\usepackage{microtype}
\usepackage{natbib}
\usepackage{url}
\usepackage[colorlinks=true,allcolors=blue!55!black]{hyperref}

\newtheorem{definition}{Definition}
\newtheorem{theorem}{Theorem}
\newtheorem{lemma}{Lemma}
\newtheorem{proposition}{Proposition}
\newtheorem{assumption}{Assumption}
\newtheorem{remark}{Remark}

\newcommand{\tespec}{\textsc{TeSpec}}
\newcommand{\qcp}{\textsc{QCP}}
\newcommand{\coins}{\textsc{Coins}}
\newcommand{\verina}{\textsc{Verina}}
\newcommand{\rocq}{Rocq}
\newcommand{\pass}{\textsf{PASS}}
\newcommand{\fail}{\textsf{FAIL}}
\newcommand{\unknown}{\textsf{UNKNOWN}}
\newcommand{\error}{\textsf{ERROR}}
\newcommand{\bindfail}{\textsf{BIND-FAIL}}
\newcommand{\true}{\mathsf{true}}
\newcommand{\false}{\mathsf{false}}
\newcommand{\Prop}{\mathsf{Prop}}
\newcommand{\Val}{\mathsf{Val}}
\newcommand{\Addr}{\mathsf{Addr}}
\newcommand{\State}{\mathsf{State}}
\newcommand{\Heap}{\mathsf{Heap}}
\newcommand{\Store}{\mathsf{Store}}
\newcommand{\Ret}{\mathsf{Ret}}
\newcommand{\Beh}{\mathsf{Beh}}
\newcommand{\dom}{\mathop{\mathrm{dom}}}
\newcommand{\sem}[1]{[\![ #1 ]\!]}
\newcommand{\sat}{\models}
\newcommand{\pto}{\mathrel{\mapsto}}
\newcommand{\exec}{\Downarrow}
\newcommand{\judge}{\mathcal{J}}
\newcommand{\bind}{\mathcal{B}}
\newcommand{\oracle}{\mathcal{O}}
\newcommand{\testsuite}{\mathcal{C}}
\newcommand{\ctx}{\Sigma}
\newcommand{\ghost}{\eta}
\newcommand{\pre}{\mathsf{Pre}}
\newcommand{\post}{\mathsf{Post}}
\newcommand{\goal}{\mathsf{Goal}}
\newcommand{\atom}{\mathsf{Atom}}
\newcommand{\canon}{\mathsf{canon}}
\newcommand{\replay}{\mathsf{replay}}

\lstdefinelanguage{QCP}{
  language=C,
  morekeywords={With,Require,Ensure,Let,Extern,Coq,Import,exists,forall},
  sensitive=true
}
\lstset{
  language=QCP,
  basicstyle=\ttfamily\small,
  keywordstyle=\color{blue!65!black}\bfseries,
  commentstyle=\color{green!35!black},
  numbers=left,
  numberstyle=\scriptsize\color{gray},
  frame=single,
  breaklines=true,
  columns=fullflexible,
  xleftmargin=1.5em,
  framexleftmargin=1em
}

\title{\textsc{TeSpec}: Instance-Based Evaluation of\\
Non-Executable Separation-Logic Specifications for C}
\author{Anonymous Authors}
\date{}

\begin{document}
\maketitle

\begin{abstract}
Formal specifications generated by humans or language models are difficult to
evaluate.  Full semantic equivalence to a reference specification is often
harder to prove than the specification-generation task itself, while ordinary
testing presumes that the property under test is executable.  Recent
instance-based approaches avoid full equivalence by specializing a declarative
specification to concrete input/output examples and proving the resulting
propositions.  Their program model, however, is predominantly functional:
behavior is a relation over values.  C contracts additionally describe mutable
memory, aliasing, ownership, recursively defined data structures, old-state
values, and ghost variables whose values do not occur in the program.

We present \tespec, a framework for evaluating non-executable
separation-logic contracts against concrete C state transitions.  A test
instance fixes the C arguments and an exact, canonical initial heap.
QCP's \texttt{With} variables are universally quantified contract parameters,
not ordinary existential witnesses.  Consequently, a supplied binding
represents the contract on that state only when it is proved unique, proved
postcondition-equivalent to every compatible binding, or belongs to a
ghost-complete finite set.  \tespec{} executes finite loops and visible callees
along the concrete path, specializes the elaborated contract with the observed
final state, and discharges the resulting verification condition using
\qcp, SMT, and a \rocq{} proof checker.  Acceptance and rejection are
bidirectional proof obligations: inability to prove either polarity yields
\unknown, never success.  Language models may propose bindings or residual
proofs, but a conclusive result requires a kernel-checked certificate.

This paper formalizes heap-transition specification evaluation, identifies four
independent notions of pre- and postcondition quality, proves conditional
soundness of the three-valued judgment, and defines a benchmark protocol that
prevents candidate specifications from controlling their own test heaps.  The
result generalizes instance-based specification evaluation from pure
input/output relations to stateful, heap-manipulating C programs without
requiring specifications to be executable.
\end{abstract}

\section{Introduction}

Formal verification can establish that an implementation satisfies a
mathematical contract, but this guarantee is only as meaningful as the
contract.  A vacuous precondition can exclude every interesting execution; a
weak postcondition can admit faulty results; and a syntactically sophisticated
contract can still formalize the wrong intent.  This problem becomes especially
visible when large language models (LLMs) generate specifications: syntax and
type checking are easy to measure, whereas semantic quality is not.

Existing evaluation methods make different compromises.  Equivalence-based
benchmarks ask for bidirectional implication between a generated and a
reference specification, as in \textsc{Clever}~\citep{thakur2025clever}.  This
provides a strong result when the proof succeeds, but proof search can dominate
the task.  Program-verification-based evaluation checks that a reference
implementation satisfies the generated contract, as in
FormalBench~\citep{lecong2025formalbench}; it detects contracts that are too
strong but requires universal verification and may not expose contracts that
are too weak.  Instance-based approaches such as
\verina~\citep{ye2026verina} and \coins~\citep{yang2026coins} specialize a
candidate proposition to trusted examples.  This substantially lowers proof
complexity and directly tests whether a specification admits desired behavior
and excludes undesired behavior.

The last approach is attractive, but a C separation-logic contract is not
merely a proposition $S(i,o)$ over immutable values.  It relates an initial
store and heap to a return value and a final store and heap.  Its spatial
assertions may describe disjoint ownership, arrays, nested structures, and
user-defined inductive predicates.  Its \texttt{With} variables may range over
arbitrary logical types and represent an abstract view of the initial heap.
They have no runtime value to record.  Moreover, obtaining a concrete final
heap may require executing loops and function calls that conventional
deductive verification handles through invariants and callee contracts.

This paper asks:

\begin{quote}
\emph{Can we evaluate a declarative C heap specification on concrete behaviors
without making the specification executable and without universally verifying
the implementation?}
\end{quote}

Our answer is \tespec.  The key object is a \emph{heap transition}
\[
  (a,\sigma_0) \longmapsto (r,\sigma_1),
\]
not a pure input/output pair.  The test fixes a binding for every logical input
before executing the program; the controlled benchmark additionally certifies
that this binding is an adequate representative of the contract's universally
quantified ghost parameters.  A modified \qcp{} executor follows the finite
concrete path, directly unfolds loops and visible callees, and carries the
separation heap through the execution.  At the terminal state it checks the
elaborated \texttt{Ensure}.  Closed logical operations are normalized,
suitable goals are discharged by \qcp/SMT, and expressive residual
propositions remain as \rocq{} obligations.

The method deliberately separates \emph{evidence} from \emph{search}.
An LLM may infer how a concrete list corresponds to a polymorphic ghost value,
or suggest a proof of a residual quantified proposition.  It cannot report a
test as passed.  A result is accepted only when a proof-producing automatic
path or the proof-assistant kernel checks the corresponding proposition.

This paper makes the following contributions:

\begin{enumerate}[leftmargin=1.4em]
  \item We define instance-based evaluation for separation-logic contracts over
        C state transitions, including arbitrary typed ghost variables and
        old-state references.
  \item We formalize the universal semantics of QCP \texttt{With} variables
        and distinguish a typed binding from an \emph{adequate
        representative}.  Pre-state-only binding prevents output leakage;
        uniqueness, observational equivalence, or ghost completeness prevents
        a single convenient abstraction from standing for incompatible ones.
  \item We define a bidirectional, three-valued judgment.  \pass{} requires a
        proof of the expected polarity, \fail{} requires a proof of its
        negation, and all inconclusive cases remain \unknown.
  \item We define trace-specialized execution over QCP's four exit modes and
        give local soundness obligations for loads, stores, finite loop
        unrolling, and visible-call inlining.  This explains precisely why one
        finite derivation needs neither a loop invariant nor a callee
        specification.
  \item We give a benchmark protocol with candidate-independent canonical
        heaps, positive and negative pre/post cases, proof provenance, and
        explicit separation of specification, binding, and proving ability.
\end{enumerate}

The paper does not claim that passing finitely many instances proves universal
equivalence.  Rather, it defines a sound evidence layer: every conclusive
instance verdict has a precise semantic meaning, and the coverage of a finite
suite is reported separately.

\section{Motivating Example}
\label{sec:example}

Consider a conventional inductive list predicate:
\[
\begin{array}{rcl}
\mathsf{sll}(p,[]) &\triangleq& p=\mathsf{null}\land emp,\\
\mathsf{sll}(p,x::xs) &\triangleq&
  \exists q.\;p\pto\{\mathsf{data}:x,\mathsf{next}:q\}
  * \mathsf{sll}(q,xs).
\end{array}
\]
The following C function increments the first element.

\begin{lstlisting}
struct node { int data; struct node *next; };

void bump_head(struct node *p)
/*@ With (x: Z) (xs: list Z)
    Require p != 0 && sll(p, cons(x, xs))
    Ensure  sll(p, cons(x + 1, xs))
*/
{
    p->data = p->data + 1;
}
\end{lstlisting}

A concrete instance binds $p$ to a root block, $x$ to $4$, and $xs$ to
$[8,15]$.  Its initial heap is
\[
b_0.\mathsf{data}=4,\quad b_0.\mathsf{next}=b_1,\quad
b_1.\mathsf{data}=8,\quad b_1.\mathsf{next}=b_2,\quad
b_2.\mathsf{data}=15,\quad b_2.\mathsf{next}=\mathsf{null}.
\]
Executing the body yields the same heap except
$b_0.\mathsf{data}=5$.  The postcondition is not an executable Boolean
predicate: it includes an inductive spatial assertion and an existential
sequence of next pointers.  Nevertheless, after fixing the two heaps and the
ghost values, it is a closed logical proposition that can be simplified and
proved.

The example exposes three evaluation hazards.

\paragraph{Candidate-controlled input.}
If the initial heap is materialized separately from each candidate's
\texttt{Require}, a candidate that accidentally describes a one-node list may
be evaluated on a different heap from a candidate that describes a three-node
list.  Cross-candidate scores would then be incomparable.  The benchmark must
freeze the initial heap independently of the candidate.

\paragraph{Post-state witness leakage.}
Suppose a binding agent sees the final heap before choosing $xs$.  It could set
$xs$ to whichever tail makes the candidate postcondition provable.  The test
would measure retrospective witness fitting, not specification quality.

\paragraph{One-sided proof search.}
If a prover fails to establish the postcondition, the proposition may be false
or the prover may simply be incomplete.  Similarly, failure to prove that a
bad transition is accepted is not evidence that the candidate rejects it.
Conclusive rejection requires a proof of the opposite polarity.

These hazards motivate the formal model below.

\section{Background and Scope}

\subsection{Separation-logic contracts}

Separation logic describes stores and heaps with local reasoning about disjoint
memory~\citep{reynolds2002separation}.  Assertions include pure propositions,
points-to atoms, separating conjunction $A*B$, and inductive predicates.
Checking whether a concrete heap satisfies symbolic-heap formulas with
inductive predicates is itself a well-studied model-checking
problem~\citep{brotherston2016model}.  Real C verification languages, however,
combine the symbolic-heap fragment with arithmetic, old-state values,
quantification, abstract logical functions, and domain-specific definitions.
The resulting contract need not be decidable or executable.

\qcp{} is a separation-logic-based C verification tool built on the Verified
Software Toolchain lineage~\citep{app2011vst,wu2026qcp}.  It uses symbolic
execution, spatial strategies, an SMT layer, and residual \rocq{} proofs.
A typical function annotation has the form
\[
\texttt{With }\bar{\alpha}\quad
\texttt{Require }P(\bar{\alpha})\quad
\texttt{Ensure }Q(\bar{\alpha}).
\]
The logical variables $\bar{\alpha}$ can range over types such as integers,
lists, nested constructors, or application-specific \rocq{} types.  They are
not C parameters.

Three details of QCP's architecture matter to our metatheory.  First, its
C-like user assertions elaborate to internal separation-logic assertions:
surface field reads carry spatial access information and may introduce
logical values.  Second, symbolic execution computes strongest postconditions
over \emph{partial statements} while tracking four exit assertions---Normal,
Break, Continue, and Return.  Memory accesses and ordinary function calls
generate spatial entailments; standard QCP uses loop invariants and callee
specifications to summarize them.  Third, QCP emits all VC statements in
\texttt{proof\_goal.v}, automatic proofs in \texttt{proof\_auto.v}, residual
proofs in \texttt{proof\_manual.v}, and a \texttt{proof\_check.v} module that
checks that every declared VC is implemented by one of the proof modules.
Its Stellis-based spatial strategies generate soundness lemmas for deduction
steps.  The current SMT component covers LIA and uninterpreted functions, but
not quantified formulas~\citep{wu2026qcp}.

These facts justify reusing QCP's assertion library, checked entailment proofs,
and proof-coverage architecture.  They do \emph{not} by themselves prove the
soundness of TeSpec's modified loop/call executor, state encoder, or
bidirectional case generator.  Section~\ref{sec:tcb} keeps those obligations
separate instead of claiming wholesale inheritance of QCP soundness.

\subsection{What is evaluated}

\tespec{} evaluates a \emph{candidate contract}.  The C implementation is a
behavior generator: a trusted implementation produces positive transitions,
while injected-bug or mutated implementations can produce negative
transitions.  The implementation is not the artifact being scored.

Our primary scope is deterministic, terminating, sequential C functions whose
relevant executions use:
\begin{itemize}[leftmargin=1.4em]
  \item scalar values, pointers, arrays, and closed structures;
  \item singly and doubly linked lists and recursive combinations of the above;
  \item finite loops, nested loops, visible calls, and finite recursion; and
  \item QCP-defined or \rocq-defined logical predicates.
\end{itemize}
Trees, unrestricted graphs, concurrency, volatile or device memory, undefined
C behavior, and opaque external effects are outside the present semantic
claim.  They may be added through explicit trusted models but are never guessed
by the evaluator.

\section{Semantic Model}
\label{sec:model}

\subsection{Concrete C states and transitions}

Let $\Val$ contain typed scalar values and pointers.  We model a pointer as a
block-offset pair $(b,o)$ rather than a raw numeric address.  A store
$\rho\in\Store$ maps local and global identifiers to values.  A finite heap
$h\in\Heap$ maps allocated block-offset pairs to typed cells.  A machine state
is $\sigma=(\rho,h)\in\State$.

For a function $f$, let $a$ be a tuple mapping its formal parameters to
concrete values.  We write
\[
  \langle f,a,\sigma_0\rangle \exec (r,\sigma_1)
\]
when the C semantics terminates safely with return value $r$ and final state
$\sigma_1$.  Fault and divergence are explicit nonterminal outcomes.  A
terminating behavior is
\[
  \beta=(f,a,\sigma_0,r,\sigma_1)\in\Beh.
\]
The benchmark stores $\sigma_0$ and $\sigma_1$ canonically; Section
\ref{sec:canonical} defines when representations denote the same heap.

\subsection{Surface assertions, elaboration, and contracts}

Let $\ctx$ be a typed logical context containing the definitions of pure
functions, spatial predicates, inductive types, and the QCP C memory model.  A
candidate may refer to definitions supplied by the task.  Candidate-supplied
extensions must be transparent definitions or positivity-checked inductive
declarations; they may not introduce axioms, parameters used as semantic
oracles, admitted lemmas, or an alternate memory semantics.

QCP source annotations are not themselves the semantic assertion language.
Let $U$ range over surface assertions and $A$ over internal
separation-logic assertions.  We write
\[
  \ctx;\Gamma_C;\Gamma_G \vdash U \Downarrow_{\!A} A
\]
for type-directed elaboration.  For example, a C-like field equality may
elaborate to an existential value, a pure equality, and a points-to or
permission atom for that field.  Elaboration also resolves C integer types,
addresses, old-state expressions, and overloaded assertion notation.

\begin{assumption}[Elaboration adequacy]
For every well-typed supported surface assertion $U$, state $\sigma$, and
ghost environment $\ghost$,
\[
  \sigma,\ghost \sat_{\mathrm{surface}} U
  \quad\Longleftrightarrow\quad
  \sigma,\ghost \sat_{\ctx} A
  \qquad
  \text{when }\ctx;\Gamma_C;\Gamma_G\vdash U\Downarrow_{\!A}A.
\]
\end{assumption}
Unsupported or ambiguous elaboration produces \error{}, never an
approximated assertion.  The QCP frontend is part of the trusted chain until
this assumption is mechanized for the TeSpec input fragment.

After elaboration, a contract is a tuple
\[
  S=(\Gamma_G,P,Q),
\]
where
\[
\begin{split}
P &: \sem{\Gamma_G}\times(A\times\State) \rightarrow \Prop,\\
Q &: \sem{\Gamma_G}\times(A\times\State)
      \times(\Ret\times\State)\rightarrow \Prop.
\end{split}
\]
Write $x=(a,\sigma_0)$ for an input state and $y=(r,\sigma_1)$ for a terminal
outcome.  The expression \texttt{e@pre} is interpreted in $\sigma_0$;
ordinary post-state memory expressions are interpreted in $\sigma_1$.

Separating conjunction has its standard heap-splitting meaning:
\[
(\rho,h),\ghost\sat A*B
\quad\Longleftrightarrow\quad
\exists h_1,h_2.\;
h=h_1\uplus h_2\land
(\rho,h_1),\ghost\sat A\land
(\rho,h_2),\ghost\sat B.
\]
Custom predicates obtain their meaning from $\ctx$, not from their names and
not from an LLM.

\subsection{Universal semantics of \texttt{With}}
\label{sec:binding}

The logical variables introduced by QCP \texttt{With} are implicit universal
parameters of the Hoare contract:
\[
  \forall\ghost\in\sem{\Gamma_G}.\;\{P(\ghost)\}\ f\ \{Q(\ghost)\}.
\]
They must not be confused with existential witnesses introduced inside an
\texttt{Ensure}.  For a concrete input $x$ and outcome $y$, define the
intrinsic candidate semantics:
\[
\begin{split}
\mathsf{Dom}_S(x)
  &\triangleq \exists\ghost.\;P(\ghost,x),\\
\mathsf{Conf}_S(x,y)
  &\triangleq \forall\ghost.\;P(\ghost,x)\rightarrow Q(\ghost,x,y),\\
\mathsf{Accept}_S(x,y)
  &\triangleq \mathsf{Dom}_S(x)\land\mathsf{Conf}_S(x,y).
\end{split}
\]
The domain uses existence because an input is admitted when it realizes at
least one abstract view.  Conformance is universal because the original
contract promises the postcondition for every \texttt{With} instantiation
whose precondition describes that input.  This definition is independent of
any binding algorithm.

A concrete execution does not reveal ghost values.  An implementation may
therefore use a pre-state binder
\[
  \bind_S : A\times\State
  \rightharpoonup \mathcal P_{\mathrm{fin}}(\sem{\Gamma_G})
\]
to propose a finite set of typed instantiations.  A single JSON binding is the
singleton case.  The binder may inspect $S$ and $\ctx$, but its runtime input
contains neither the return value, final state, nor case polarity.

\begin{definition}[Pre-state-only binding]
A binder is pre-state-only when any two executions sharing $x=(a,\sigma_0)$
produce exactly the same binding result, independently of their terminal
outcomes and labels.
\end{definition}

\begin{lemma}[Output noninterference]
A pre-state-only binder assigns the same ghost candidates to a correct
execution and to every mutant execution replayed from the same $x$.
\end{lemma}
\begin{proof}
Immediate from the domain of $\bind_S$.
\end{proof}

Pre-state-only binding prevents retrospective witness fitting, but is not
sufficient for semantic adequacy.  The same heap can satisfy $P$ under several
abstract views whose postconditions differ.  Define observational equivalence
at $x$ by
\[
  \ghost \equiv_{S,x} \ghost'
  \quad\Longleftrightarrow\quad
  \forall y.\;Q(\ghost,x,y)\leftrightarrow Q(\ghost',x,y).
\]

\begin{definition}[Adequate representative]
A ghost value $\widehat\ghost$ adequately represents $S$ at $x$ when
\[
  P(\widehat\ghost,x)
  \quad\land\quad
  \forall\ghost'.\;
    P(\ghost',x)\rightarrow
    \ghost'\equiv_{S,x}\widehat\ghost.
\]
\end{definition}

\begin{definition}[Ghost-complete finite binding]
A finite set $E$ is ghost-complete for $(S,x)$ when the following two checked
equivalences hold for every outcome $y$:
\[
\begin{split}
\mathsf{Dom}_S(x)
  &\leftrightarrow \bigvee_{\ghost\in E}P(\ghost,x),\\
\mathsf{Conf}_S(x,y)
  &\leftrightarrow
    \bigwedge_{\ghost\in E}\bigl(P(\ghost,x)\rightarrow Q(\ghost,x,y)\bigr).
\end{split}
\]
\end{definition}

Uniqueness of a satisfying binding implies adequate representation; adequate
representation implies that the universal postcondition can be checked using
one representative.  More generally, ghost completeness permits a finite
conjunction.  These facts require proof in the task context; they are not
inferred from a binder's confidence score.

\paragraph{Two benchmark tracks.}
In the \emph{controlled-interface track}, the task fixes $\Gamma_G$ and stores
an oracle binding together with a proof of uniqueness, adequate
representation, or ghost completeness.  This is the primary specification
score.  In the \emph{open-interface track}, a candidate may introduce
arbitrary \texttt{With} types and a human or agent may supply typed values.
Without an adequacy certificate, the result is explicitly an
\emph{indexed-instance judgment} for those values, not a judgment of the
whole candidate contract.  Failure to find a binding never proves
$\neg\mathsf{Dom}_S(x)$.

\subsection{Candidate-independent heaps}
\label{sec:fixedheap}

The precondition of a separation-logic contract can be used constructively to
materialize a heap.  This is convenient for interactive testing, but unsafe as
the basis of a comparative benchmark.

\begin{definition}[Frozen pre-state]
A benchmark case has a frozen pre-state when its canonical $\sigma_0$ is
constructed and committed before a candidate specification is evaluated, and
the candidate cannot alter $\sigma_0$ through its own precondition,
definitions, or binding choices.
\end{definition}

The benchmark may initially create $\sigma_0$ from a trusted reference
contract and reference binding.  Once canonicalized, however, all candidates
are evaluated by replaying the same state.  Candidate $P$ is checked
\emph{against} that state; it does not construct it.

\begin{remark}[Interactive materialization]
The existing contract-testing interface
\[
  \texttt{candidate spec + implementation + binds}
\]
may materialize a heap from the candidate \texttt{Require}.  Such a result is
useful evidence that the contract is internally consistent with one execution,
but it is not a candidate-comparable benchmark score.  We call this
\emph{self-materialized mode}; the primary benchmark uses \emph{frozen-replay
mode}.
\end{remark}

\subsection{Canonical heaps}
\label{sec:canonical}

Raw addresses are process-specific and should not appear in a benchmark
oracle.  A canonical heap assigns stable block identifiers and records:
\[
  (\text{block kind},\text{size},\text{typed cells},\text{root references},
   \text{alias relation}).
\]
A heap isomorphism $\mu$ is a bijection on block identifiers preserving block
kind, size, offsets, scalar values, root reachability, and pointer edges.
Pointers are renamed as $(b,o)\mapsto(\mu(b),o)$.

\begin{definition}[Address-parametric fragment]
A contract is address-parametric if its truth is invariant under heap
isomorphism.  In particular, it may use pointer equality, offsets within the
same block, and field addresses, but does not depend on the numeric
representation or total ordering of unrelated pointers.
\end{definition}

\begin{proposition}[Canonicalization invariance]
For an address-parametric assertion $A$, heap isomorphism $\mu$, and
consistently renamed store and ghost pointers,
\[
  (\rho,h),\ghost\sat A
  \quad\Longleftrightarrow\quad
  (\mu(\rho),\mu(h)),\mu(\ghost)\sat A.
\]
\end{proposition}
\begin{proof}[Proof sketch]
By structural induction on $A$.  Pure scalar cases are unchanged; pointer
equality and block offsets are preserved by $\mu$; points-to atoms are
preserved by cell and edge preservation; separating conjunction transports
disjoint heap partitions through the bijection; and inductive predicates
follow by induction on their derivations.
\end{proof}

Programs or contracts that cast pointers to integers, compare unrelated pointer
representations, or depend on allocator-chosen addresses are excluded from this
proposition.  They can still be run with a fixed concrete layout, but their
result is layout-specific and must be labeled as such.

\subsection{Exact-state encoding}
\label{sec:exact}

Let $\mathsf{Enc}(\sigma)$ be an internal assertion encoding every allocated
block, typed cell, root, and relevant store value of the finite canonical
state.  ``Exact'' means that its models form precisely the heap-isomorphism
class of $\sigma$; it has no unconstrained spatial frame.

\begin{assumption}[Encoding adequacy]
For every supported, address-parametric internal assertion $A$,
\[
  \sigma\sat_{\ctx} A
  \quad\Longleftrightarrow\quad
  \mathsf{Eval}_{\ctx}(\mathsf{Enc}(\sigma),A),
\]
where $\mathsf{Eval}_{\ctx}$ is the semantic proposition emitted to \rocq{}.
\end{assumption}

The right-hand side may be implemented using QCP assertion entailment plus an
exactness lemma, but it must be a proposition in the checked logic.  A
successful proof of the positive entailment establishes satisfaction.  Failure
of that proof does \emph{not} establish nonsatisfaction.  A negative verdict
requires a proof of
\[
  \neg\mathsf{Eval}_{\ctx}(\mathsf{Enc}(\sigma),A)
\]
or a checked countermodel together with a completeness theorem for the
countermodel procedure.

A partial assertion $\mathsf{Enc}_L(\sigma)*F$ with symbolic frame $F$ denotes
a set of heap completions, not one concrete state.  It is admissible for the
primary benchmark only when the candidate assertion is proved invariant under
changes outside $L$, or when the generated proposition explicitly quantifies
over all completions.  Otherwise it provides execution information but not an
exact specification-evaluation instance.

\section{What It Means for a Specification to Be Correct}
\label{sec:quality}

\subsection{Reference semantics}

Let $R$ be the intended contract semantics.  It induces a set of valid inputs
$I_R\subseteq A\times\State$ and, for each valid input, a set of allowed
terminal outcomes:
\[
  T_R(x)\subseteq \Ret\times\State,\qquad x=(a,\sigma_0).
\]
$R$ may be represented by a trusted formal contract, expert-labeled cases, or
an ensemble of reference implementations plus checked postconditions.  The
evaluator only uses cases whose labels are certified.

A candidate $S$ intrinsically induces:
\[
\begin{split}
I_S
  &=\{x\mid \mathsf{Dom}_S(x)\},\\
T_S(x)
  &=\{y\mid\mathsf{Conf}_S(x,y)\}.
\end{split}
\]
We intentionally do not conjoin $\mathsf{Dom}_S$ into $T_S$.  Precondition and
postcondition quality are scored independently; otherwise a candidate could
make every postcondition test vacuous by setting $P=\false$.

\subsection{Four dimensions}

We use set inclusion to avoid ambiguity in the words ``sound'' and
``complete,'' whose direction differs across prior work.

\begin{definition}[Input coverage]
$S$ has input coverage with respect to $R$ when
\[
  I_R\subseteq I_S.
\]
It does not incorrectly exclude intended inputs.
\end{definition}

\begin{definition}[Input precision]
$S$ has input precision with respect to $R$ when
\[
  I_S\subseteq I_R.
\]
It does not admit inputs forbidden by the intended contract.
\end{definition}

\begin{definition}[Transition coverage]
$S$ has transition coverage with respect to $R$ when
\[
  \forall x\in I_R.\;T_R(x)\subseteq T_S(x).
\]
It accepts every intended outcome on a valid input.
\end{definition}

\begin{definition}[Transition precision]
$S$ has transition precision with respect to $R$ when
\[
  \forall x\in I_R.\;T_S(x)\subseteq T_R(x).
\]
It rejects unintended return values and heap effects on valid inputs.
\end{definition}

\begin{theorem}[Extensional correctness]
A candidate has all four properties exactly when
\[
I_S=I_R
\quad\text{and}\quad
\forall x\in I_R.\;T_S(x)=T_R(x).
\]
\end{theorem}
\begin{proof}
Equality is equivalent to mutual inclusion for both the input set and each
outcome set.
\end{proof}

Full extensional correctness is generally as difficult as semantic
equivalence.  A benchmark approximates the four inclusions using a finite set
of certified instances.

\subsection{Certified test cases}

A case has one of four kinds:
\[
\begin{array}{lll}
\mathsf{Pre}^{+}:  & x                     & x\in I_R,\\
\mathsf{Pre}^{-}:  & x                     & x\notin I_R,\\
\mathsf{Post}^{+}: & (x,r,\sigma_1)        &
                        x\in I_R\land(r,\sigma_1)\in T_R(x),\\
\mathsf{Post}^{-}: & (x,r,\sigma_1)        &
                        x\in I_R\land(r,\sigma_1)\notin T_R(x).
\end{array}
\]
For postcondition cases, positive and negative outcomes sharing $x$ also share
the same pre-state binding evidence.  The expected closed proposition is:
\[
\goal_S(c)=
\begin{cases}
\mathsf{Dom}_S(x)
  & c=\mathsf{Pre}^{+}(x),\\
\neg\mathsf{Dom}_S(x)
  & c=\mathsf{Pre}^{-}(x),\\
\mathsf{Conf}_S(x,y)
  & c=\mathsf{Post}^{+}(x,y),\\
\neg\mathsf{Conf}_S(x,y)
  & c=\mathsf{Post}^{-}(x,y).
\end{cases}
\]
Post goals deliberately omit $\mathsf{Dom}_S$: input and transition quality
are scored independently.  If $P$ is unsatisfiable at $x$, then
$\mathsf{Conf}_S(x,y)$ is true for every $y$; such a candidate accepts positive
outcomes but cannot reject any negative outcome, while separately failing the
positive-pre case.  It therefore cannot pass a suite containing both input and
transition polarities by choosing $\false$ as its precondition.

Under an adequate representative $\widehat\ghost$ for an admitted input, the
postcondition goal reduces to $Q(\widehat\ghost,x,y)$ or its negation.  Under a
ghost-complete set it reduces to the corresponding finite conjunction.
Without such evidence, replacing the quantified goals by one supplied binding
is unsound.

Thus positive cases test coverage and negative cases test precision.  A failed
binder call is an evaluation failure, not evidence for a negative pre-case.

\section{A Bidirectional Three-Valued Judgment}
\label{sec:judgment}

Let $\ctx\vdash\phi$ mean that a proof object for the closed proposition
$\phi$ is accepted in the trusted logical context.  Automated tactics and LLMs
may search for the object, but do not change this relation.

\begin{definition}[Instance judgment]
For a well-typed closed goal $\phi=\goal_S(c)$,
\[
\judge_{\ctx}(\phi)=
\begin{cases}
\pass & \text{if a checked proof of }\ctx\vdash\phi\text{ is found},\\
\fail & \text{if a checked proof of }\ctx\vdash\neg\phi\text{ is found},\\
\unknown & \text{otherwise}.
\end{cases}
\]
Parse, dependency, type, binding, replay, and environment failures are reported
separately as \error{} or \bindfail.
\end{definition}

The second line is essential.  On a negative post-state, the target is
$\neg\mathsf{Conf}_S(x,y)$.  A \pass{} therefore requires a proof that
the candidate rejects the bad transition, while a \fail{} requires a proof of
the opposite case goal.  Merely failing to establish a QCP entailment, timing
out in SMT, or receiving a symbolic-execution diagnostic establishes neither
polarity.

\begin{theorem}[Verdict exclusivity]
If $\ctx$ is consistent, no closed goal can receive both a valid \pass{}
certificate and a valid \fail{} certificate.
\end{theorem}
\begin{proof}
Such certificates would prove both $\phi$ and $\neg\phi$, contradicting
consistency.
\end{proof}

\begin{theorem}[Instance-verdict soundness]
Assume the proof checker is sound with respect to the semantics of $\ctx$.
Then:
\[
\judge_{\ctx}(\goal_S(c))=\pass
\Longrightarrow \goal_S(c)\text{ is true},
\]
and
\[
\judge_{\ctx}(\goal_S(c))=\fail
\Longrightarrow \goal_S(c)\text{ is false}.
\]
\end{theorem}
\begin{proof}
The first implication follows from soundness of a checked proof of $\goal_S(c)$.
For the second, a checked proof of $\neg\goal_S(c)$ entails that
$\goal_S(c)$ is false.
\end{proof}

\begin{proposition}[Monotonicity in proof budget]
Increasing the search budget can replace \unknown{} with \pass{} or \fail{},
but cannot soundly reverse a previously certified \pass{} or \fail{}.
\end{proposition}

An inconsistent imported context invalidates this guarantee.  Therefore
candidate-generated axioms, \texttt{Admitted}, proof escape hatches, and
untrusted opaque declarations are rejected.  Shared domain definitions are
part of the benchmark's trusted context and are versioned with the task.
In particular, a candidate-controlled \texttt{Extern Coq} proposition cannot
be accepted as an unexplained axiom; it must resolve to a whitelisted task
definition or a transparent candidate definition.

\subsection{Aggregating finite evidence}

For a candidate $S$ and finite case set $\testsuite$, define:
\[
\begin{split}
\mathsf{Resolved}(S,\testsuite)
  &= |\{c\in\testsuite\mid\judge(\goal_S(c))\neq\unknown\}|,\\
\mathsf{Passed}(S,\testsuite)
  &= |\{c\in\testsuite\mid\judge(\goal_S(c))=\pass\}|,\\
\mathsf{Failed}(S,\testsuite)
  &= |\{c\in\testsuite\mid\judge(\goal_S(c))=\fail\}|.
\end{split}
\]
Scores are reported separately for all four case kinds.  A strict task-level
success requires every case to be \pass; \unknown{} is never silently counted
as success.

When comparison with systems using incomplete proof search requires a range,
we report:
\[
\mathsf{LB}=\frac{\mathsf{Passed}}{|\testsuite|},
\qquad
\mathsf{UB}=\frac{\mathsf{Passed}+\mathsf{Unknown}}{|\testsuite|}.
\]
These bounds describe evaluator uncertainty, not a probabilistic confidence
interval and not a universal correctness guarantee.

\section{Trace-Specialized Specification Evaluation}
\label{sec:trace}

\subsection{Why loops need no invariants here}

A loop invariant summarizes arbitrarily many iterations in a universal proof.
A concrete test observes one finite iteration count.  If each branch guard and
memory read needed by that execution is concrete, the evaluator can repeatedly
execute the body until the guard becomes false.  This establishes one trace,
not the loop's behavior on all states.

The same distinction applies to calls.  A callee contract summarizes all
callee executions.  For one call with a visible body and concrete actual
arguments, \tespec{} enters a fresh frame, executes the body, carries the
shared heap into and out of the call, and returns the observed value.  No
callee logical binding is requested.  The caller's actual arguments determine
the callee's runtime parameters.

\subsection{Execution closure and evaluation closure}

Let $\tau$ be a concrete trace,
$\mathsf{Foot}_{\mathrm{exec}}(\tau)$ contain every location whose content or
validity can affect that trace, and $\mathsf{Foot}_{\mathrm{obs}}(S,c)$ contain
every location that the elaborated candidate goal may observe.

\begin{definition}[Execution closure]
A case is execution-closed for $f$ if:
\begin{enumerate}[leftmargin=1.4em]
  \item every C argument is concrete;
  \item every location in $\mathsf{Foot}_{\mathrm{exec}}(\tau)$ has a
        concrete, well-typed cell and allocation status;
  \item every encountered guard reduces to a unique concrete Boolean;
  \item every called body is visible and execution-closed, or has an
        explicitly trusted external transition model; and
  \item the trace terminates within the loop, call-depth, and recursion bounds.
\end{enumerate}
\end{definition}

\begin{definition}[Evaluation closure]
A case is evaluation-closed for candidate $S$ when it is execution-closed and
the replay fixes every location in
\[
  \mathsf{Foot}_{\mathrm{exec}}(\tau)
  \cup\mathsf{Foot}_{\mathrm{obs}}(S,c),
\]
or provides a checked frame-invariance theorem for the omitted region.
\end{definition}

Execution closure is enough to choose the concrete program path.  It is not
enough to evaluate an arbitrary candidate specification: a candidate may
mention memory never read by the program.  The primary benchmark therefore
uses the exact encoding of Section~\ref{sec:exact}.  Partial instantiation is
permitted as an optimization only after observation independence is proved;
otherwise the result is universal over symbolic completions and must be
labeled accordingly.

\subsection{Specialization pipeline}

For a candidate $S$, case $c$, and frozen pre-state, \tespec{} performs:

\begin{enumerate}[leftmargin=1.4em]
  \item \textbf{Bind.}  Type-check the C arguments, ghost types, and logical
        values.  In the open track, invoke $\bind_S(a,\sigma_0)$ before
        execution.
  \item \textbf{Replay.}  Decode the canonical pre-state into fresh concrete
        blocks while preserving roots, offsets, and aliases.
  \item \textbf{Execute.}  Symbolically execute the C syntax with a concrete
        store and touched heap.  Directly unroll finite loops and inline visible
        callees.  Record safety conditions and the terminal state.
  \item \textbf{Specialize.}  Substitute $a,\ghost,\sigma_0,r,\sigma_1$ into
        the selected $P$ or $Q$.  Preserve the original logical definition of
        the contract; do not replace it by an executable surrogate.
  \item \textbf{Normalize.}  Apply semantics-preserving reductions to closed
        arithmetic, list operations, constructor equalities, finite structural
        unfoldings, and concrete heap atoms.
  \item \textbf{Discharge both polarities.}  Try $\goal_S(c)$ and, if needed,
        $\neg\goal_S(c)$ through proof-producing QCP/SMT automation.  Emit any
        unresolved proposition as a residual \rocq{} goal.
  \item \textbf{Check and classify.}  Validate proof artifacts, record their
        provenance, and return \pass, \fail, or \unknown.
\end{enumerate}

Intermediate loop invariants and callee contracts are not used as execution
summaries.  Internal assertions can be ignored when the benchmark target is
only the selected function contract, but memory-safety and definedness
conditions needed to justify the observed C execution cannot be skipped.

\subsection{Generic normalization}

Specialization is valuable only if it avoids expression blow-up without
embedding benchmark-specific rewrites.  We require each normalization rule
$N$ to satisfy:
\[
  \ctx\vdash \phi\leftrightarrow N(\phi).
\]
The intended rule families operate on syntax classes rather than program or
predicate names:

\begin{itemize}[leftmargin=1.4em]
  \item evaluation of closed integer, Boolean, and supported floating
        expressions;
  \item reduction of closed constructor matches, list length, indexing, and
        slicing operations;
  \item first-order unification of constructor equalities;
  \item guarded unfolding of transparent recursive definitions on known
        constructors;
  \item cancellation and reconstruction of concrete spatial atoms; and
  \item bounded enumeration only when a finite domain follows from the closed
        proposition and remains below a global resource limit.
\end{itemize}

Rules do not inspect function names, case identifiers, or concrete benchmark
constants.  Unsupported operations remain residual; they are not approximated
by $\true$.

\subsection{Concrete specialization of QCP symbolic execution}

QCP symbolically executes partial statements while maintaining Normal, Break,
Continue, and Return assertions.  TeSpec retains these exit modes but
specializes the active assertion to an exact concrete state.  Write
\[
  \ctx\vdash
  \langle ps,A\rangle
  \Longrightarrow_k
  \langle\kappa,A'\rangle;\mathcal O,
  \qquad
  \kappa\in\{\mathsf N,\mathsf B,\mathsf C,\mathsf R(v)\},
\]
when partial statement $ps$, with fuel $k$, produces exit assertion $A'$ and
side obligations $\mathcal O$.  Every conclusive derivation requires a checked
proof of every obligation.

For a load from location $\ell$, the spatial obligation has the QCP form
\[
  A\vdash\exists v,F.\;\ell\pto v * F.
\]
Concrete mode additionally requires $v$ to reduce uniquely to a ground typed
value.  A store replaces $\ell\pto v*F$ by $\ell\pto v'*F$ after proving the
address, permission, and C definedness conditions.  If an address, loaded
value, or guard remains ambiguous, execution returns \unknown{} rather than
choosing a model.

The loop rule is a finite derivation, not an induction invariant.  Omitting
assertions for brevity, its defining cases are:
\[
\begin{array}{ll}
\sem{e}_\sigma=\mathsf{false}
  &\Rightarrow
    \langle\mathbf{while}\ e\ s,\sigma\rangle
      \Downarrow_k(\mathsf N,\sigma),\\
\sem{e}_\sigma=\mathsf{true},\
  \langle s,\sigma\rangle\Downarrow(\mathsf B,\sigma')
  &\Rightarrow
    \langle\mathbf{while}\ e\ s,\sigma\rangle
      \Downarrow_k(\mathsf N,\sigma'),\\
\sem{e}_\sigma=\mathsf{true},\
  \langle s,\sigma\rangle\Downarrow(\mathsf R(v),\sigma')
  &\Rightarrow
    \langle\mathbf{while}\ e\ s,\sigma\rangle
      \Downarrow_k(\mathsf R(v),\sigma'),\\
\sem{e}_\sigma=\mathsf{true},\
  \langle s,\sigma\rangle\Downarrow(\kappa,\sigma'),\
  \kappa\in\{\mathsf N,\mathsf C\},\
  \langle\mathbf{while}\ e\ s,\sigma'\rangle
      \Downarrow_{k-1}(\kappa',\sigma'')
  &\Rightarrow
    \langle\mathbf{while}\ e\ s,\sigma\rangle
      \Downarrow_k(\kappa',\sigma'').
\end{array}
\]
Fuel exhaustion is \unknown{}.  A visible call evaluates actual arguments,
allocates a fresh local frame, derives the callee body, propagates its heap and
return value, and then discards the local frame.  It does not instantiate the
callee's \texttt{With} variables because no callee contract is used.

\begin{theorem}[Finite-loop trace soundness]
Assume sound expression evaluation and sound local load/store rules.  If all
side obligations in a finite loop derivation are valid, the derived exit and
state agree with the supported C operational semantics.
\end{theorem}
\begin{proof}[Proof sketch]
By induction on the number of true guard evaluations.  The base case is the
false-guard rule.  The step uses body soundness; Break and Return are terminal
cases, whereas Normal and Continue apply the induction hypothesis to the
strictly smaller remaining derivation.
\end{proof}

\begin{theorem}[Visible-call trace soundness]
Assume type-correct argument evaluation, fresh local-frame allocation, and a
sound derivation of the visible callee body.  The call rule returns the same
value and shared-heap effect as the supported C operational semantics.
\end{theorem}
\begin{proof}[Proof sketch]
By induction on call depth, using frame freshness to show that discarding
callee locals preserves the caller store and shared heap.
\end{proof}

Let
\[
\mathsf{ExecVC}(f,a,\sigma_0,r,\sigma_1,\tau)
\]
conjoin the local side obligations and the derived terminal state.  The two
theorems above reduce global trace soundness to the remaining singleton
statement rules.  A publication artifact must either mechanize this induction
for the supported QCP partial-statement semantics or explicitly retain it as a
TCB assumption.  QCP's ordinary proof-check file verifies all \emph{generated}
VCs, but does not establish that a modified generator emitted every required
VC.

\begin{lemma}[Specialization preservation]
Let $\theta$ substitute well-typed concrete arguments, ghost values, and
states into $P$ or $Q$.  If normalization produces $N(\theta(\phi))$ using
equivalence-preserving rules, then:
\[
\ctx\sat \theta(\phi)
\quad\Longleftrightarrow\quad
\ctx\sat N(\theta(\phi)).
\]
\end{lemma}
\begin{proof}
By composition of substitution soundness and the equivalence theorem for every
normalization step.
\end{proof}

\begin{theorem}[Terminal-contract adequacy]
Assume execution-rule soundness, encoding and elaboration adequacy,
specialization preservation, and a checked
proof of
\[
\mathsf{ExecVC}(f,a,\sigma_0,r,\sigma_1,\tau)
\land \mathsf{Conf}_S(x,y).
\]
Then the observed C execution terminates at $(r,\sigma_1)$ and satisfies the
candidate contract on $x$ and $y$.
\end{theorem}
\begin{proof}
The local statement, loop, and call soundness theorems establish the C
execution.  Encoding, elaboration, and specialization preservation connect the
checked closed proposition to the original contract semantics.
\end{proof}

The theorem does not establish that all executions of $f$ satisfy $Q$, nor that
$Q$ matches user intent outside the evaluated case set.

\subsection{What is inherited from QCP}

QCP's \texttt{proof\_goal}/\texttt{proof\_auto}/\texttt{proof\_manual}/
\texttt{proof\_check} split provides a strong \emph{proof-coverage} property:
if the final module check succeeds without proof escapes, each VC field listed
in the goal interface is implemented by a kernel-accepted proof.  Automatic
spatial strategies are likewise safe when their generated soundness lemmas and
the resulting proof are checked against the original VC.  This architecture
does not validate omissions or mistranslations by a changed frontend or VC
generator.

\begin{center}
\small
\begin{tabular}{p{.22\linewidth}p{.27\linewidth}p{.39\linewidth}}
\toprule
Component & QCP evidence & TeSpec obligation/status\\
\midrule
Internal memory and separation logic
  & Rocq definitions and tactics
  & Reused; consistency of task imports remains assumed.\\
Surface assertion frontend
  & QCP implementation and examples
  & Elaboration adequacy must be established for the accepted fragment.\\
Spatial strategies
  & Stellis deduction steps generate soundness lemmas
  & Reused only through checked proofs of the original goals.\\
SMT discharge
  & LIA+UF solver contributes automatic VC proofs
  & Counted as conclusive only when represented in the checked proof module.\\
Proof aggregation
  & Module Type covers every generated goal field
  & Reused; final \texttt{proof\_check.v} compilation is mandatory.\\
Finite loops and visible calls
  & Standard QCP instead uses invariants and callee contracts
  & Modified rules; require the trace-soundness theorems above.\\
Exact state and ghost binding
  & Not a standard QCP responsibility
  & New encoding, elaboration, and binding-adequacy obligations.\\
Negative polarity
  & Ordinary verification reports failure to derive
  & New dual goal; solver failure is never a \fail{} certificate.\\
\bottomrule
\end{tabular}
\end{center}

\section{Positive and Negative Heap Behaviors}
\label{sec:negative}

Positive cases alone principally expose specifications that are too strong.
Transition precision requires negative outcomes.

\subsection{Positive behaviors}

A positive input is certified by proving the reference precondition on its
frozen state.  A positive transition is obtained by executing a trusted
reference implementation and proving the reference postcondition on the
observed terminal state.  The candidate is then evaluated on exactly those
states.

\subsection{Negative inputs}

A negative input consists of concrete arguments and a pre-state for which the
reference precondition is provably false.  No target execution is required:
the candidate precondition is specialized directly to the frozen state.
Examples include null roots, invalid length relationships, broken aliasing
requirements, insufficient array regions, and malformed list segments.

\subsection{Negative outcomes}

A negative postcondition case reuses a certified positive input and produces
an outcome that provably violates the reference postcondition.  Preferably the
outcome is feasible: it is obtained by executing a compilable mutant of the
reference program.  Heap-aware mutant families include:

\begin{itemize}[leftmargin=1.4em]
  \item wrong array index, boundary, or field;
  \item omitted, duplicated, or reordered writes;
  \item incorrect \texttt{next}/\texttt{prev}, head, or tail update;
  \item lost nodes, extra nodes, or changed aliasing;
  \item incorrect return value with an otherwise correct heap;
  \item correct return value with an incorrect heap effect;
  \item off-by-one loop iteration or early exit;
  \item replacing an old-state value by its post-state value; and
  \item removing ownership, framing, or disjointness constraints.
\end{itemize}

Arbitrary corrupted heaps can also identify logical weakness, but feasible
mutant outcomes better approximate bugs that an implementation can exhibit.
Every negative label must itself be proved against $R$; a disagreement with
one implementation is not sufficient.

\subsection{Adversarial expansion and reduction}

VeriScale shows that adversarial implementations can expose weaknesses hidden
by small test suites and that a compact discriminative subset can be recovered
after expansion~\citep{bai2026veriscale}.  The same idea lifts to heaps:
generate candidate specifications and program mutants, retain a case whenever
it separates at least one candidate from the reference, then solve a set-cover
problem over the candidate/mutant kill matrix.  This optimization changes suite
size, not case semantics.  Held-out projects and predicate families are
required to prevent the reducer from overfitting one corpus.

\section{Agent Assistance Without an Agent Judge}
\label{sec:agent}

\tespec{} separates four possible agent roles:

\begin{description}[leftmargin=1.45in,style=nextline]
  \item[Specification agent] Generates the candidate $P$ and $Q$.
  \item[Binding agent] Maps a frozen pre-state to typed \texttt{With} values in
        the open-interface track.
  \item[Proof agent] Proposes a proof term or tactic script for a residual goal.
  \item[Red-team agent] Proposes invalid inputs, heap mutants, or implementation
        mutants; the reference oracle must certify their labels.
\end{description}

Only the first role is the object being evaluated in a specification-generation
study.  The remaining roles are evaluator components and must be measured
separately.

\begin{theorem}[Agent containment]
Suppose:
\begin{enumerate}[leftmargin=1.4em]
  \item bindings are type-checked and produced through the pre-state-only API;
  \item every whole-contract verdict uses a checked adequate-representative or
        ghost-completeness certificate;
  \item candidate and proof files cannot add axioms or modify the trusted
        context;
  \item automatic and agent-generated proofs are accepted only after checking;
        and
  \item test labels are certified independently of the proposing agent.
\end{enumerate}
Then no text emitted by an agent can directly create a false conclusive
verdict.
\end{theorem}
\begin{proof}[Proof sketch]
Bindings cannot depend on the output, and their adequacy proof connects the
finite instantiated goal to the universal \texttt{With} semantics.  Candidate
text defines the proposition but cannot prove it by extending the logic.  A
conclusive verdict requires a checked proof of one polarity.  Proposed cases
affect the suite only after independent label certification.  Thus agent
output can cause failure, rejection, or incompleteness, but not bypass the
trusted judgment.
\end{proof}

Without Condition~2, agent containment applies only to the indexed proposition
that was actually instantiated.  It does not lift that proposition to the
whole contract when several incompatible bindings are possible.

\section{Benchmark Definition}
\label{sec:benchmark}

\subsection{Task package}

A task package contains:

\begin{verbatim}
id, natural_language_description
C source and target function
trusted logical context and allowed candidate interface
reference contract or certified oracle
positive pre-states and transitions
negative pre-states and transitions
canonical heaps, roots, aliases, and type layouts
case-local QCP strategies and Rocq definitions
feature labels and provenance
\end{verbatim}

Dependencies are task-local and content-addressed.  A candidate cannot shadow a
shared module.  If a specification imports a \rocq{} definition, that
definition is part of the task semantics.  QCP strategies are automation
hints: they may guide unfolding or heap cancellation but cannot silently
replace the predicate definition.

\subsection{Candidate protocol}

The controlled track fixes:
\begin{itemize}[leftmargin=1.4em]
  \item the C function signature;
  \item the \texttt{With} schema and logical types;
  \item the names and signatures of available logical definitions; and
  \item the files that the candidate may modify.
\end{itemize}
This mirrors the common practice of fixing a Lean function/specification
signature in verifiable-code benchmarks.  The open track permits new
transparent definitions and arbitrary \texttt{With} types, but reports syntax,
dependency resolution, binding, and semantic judgment as distinct stages.

\subsection{Evidence record}

For each candidate and case, the evaluator stores:
\[
\begin{array}{ll}
\text{identity} & \text{task, candidate, case, implementation/mutant hash},\\
\text{state} & \canon(\sigma_0),\canon(\sigma_1),a,r,E,
                  \text{ binding-adequacy certificate},\\
\text{logic} & \goal_S(c),\text{normalized goal, imported context hashes},\\
\text{execution} & \text{trace, safety VCs, limits, terminal-state hash},\\
\text{proof} & \text{polarity, source, checker version, certificate hash},\\
\text{verdict} & \pass,\fail,\unknown,\error,\text{ or }\bindfail.
\end{array}
\]
Automatic and residual proofs remain visibly separate.  A manually or
agent-proved residual can resolve an \unknown{}, but is not relabeled as SMT
automation.

\subsection{Metrics}

The benchmark reports:
\begin{itemize}[leftmargin=1.4em]
  \item parse, dependency, type-check, bind, replay, execution-closure, and
        evaluation-closure rates;
  \item positive-pre, negative-pre, positive-post, and negative-post pass rates;
  \item strict all-cases task success;
  \item \unknown{} rate and lower/upper evidence bounds;
  \item heap-mutant rejection score;
  \item automatic versus residual proof provenance;
  \item agent cost and proof-search budget; and
  \item results stratified by arrays, structures, singly/doubly linked lists,
        composition, calls, loops, quantifiers, and custom predicates.
\end{itemize}

The primary specification score uses oracle bindings in the controlled track.
End-to-end results with agent-generated bindings are reported separately.
Cases without binding adequacy are reported as indexed-instance results and
are excluded from the whole-contract score.

\section{Implementation Architecture}
\label{sec:implementation}

The \tespec{} prototype is built around a modified \qcp{} symbolic executor.
Its public input is a QCP-annotated C implementation, local strategy and
\rocq{} dependencies, a target function, and one or more JSON binding records.
The analyzer discovers C parameters, value-level \texttt{With} variables, and
type-level \texttt{With} parameters.  Bindings support scalars, typed nested
lists, arbitrary constructor trees, and a lossless raw QCP term for
application-specific types.

The concrete mode modifies standard deductive execution in two places.  First,
loops are driven by their concrete guards until exit rather than summarized by
an invariant.  Second, visible callees are executed in fresh frames with
call-site arguments and the shared heap rather than summarized by callee
contracts.  Both mechanisms have explicit bounds whose exhaustion yields
\unknown.  Recursive QCP \texttt{Let} predicates are unfolded by typed
substitution and constructor unification; \texttt{Extern Coq} spatial
predicates require explicit strategies or remain residual.

The backend generates the ordinary QCP/\rocq{} artifact split: goals,
proof-producing automatic steps, residual manual lemmas, and a final check
file.  Closed arithmetic and structural formulas are delegated to the existing
SMT layer.  Residual proof files are checked for proof escapes and compiled by
the \rocq{} kernel.  Quantified or domain-specific propositions are not assumed
to become easy merely because program values are concrete; if normalization
does not expose a supported LIA+UF goal, they remain residual.

\subsection{Current prototype versus benchmark semantics}

The current executable realizes the following path:
\[
\begin{gathered}
\texttt{candidate spec + implementation + binds}\\[-1mm]
\downarrow\\[-1mm]
\texttt{self-materialized pre-state}
\;\longrightarrow\;
\texttt{concrete execution}
\;\longrightarrow\;
\texttt{terminal Ensure VC}.
\end{gathered}
\]
This is already useful for testing one supplied contract and for constructing
reference behaviors.  The formal benchmark defined in this paper additionally
requires:

\begin{enumerate}[leftmargin=1.4em]
  \item exporting the materialized reference heap to a canonical case record;
  \item replaying that frozen heap without allowing candidate $P$ to modify it;
  \item proving exact-state encoding and surface-elaboration adequacy;
  \item a precondition-only evaluation path for negative inputs;
  \item explicit positive/negative post labels and shared binding records;
  \item binding uniqueness, representative adequacy, or ghost completeness;
  \item generation and checking of both $\goal$ and $\neg\goal$; and
  \item a proof certificate or checked countermodel for every \fail{} verdict.
\end{enumerate}

Until these items are implemented, prototype \pass{} results should be
described as \emph{binding-specialized contract consistency}, not as a complete
cross-candidate specification benchmark.  Making this boundary explicit
prevents the paper's semantic claim from exceeding the artifact.
In particular, current diagnostic strings such as an inconsistent symbolic
obligation or ``cannot derive'' are not proofs of the opposite case goal.
Publication-mode classification must downgrade them to \unknown{} unless the
dual \rocq{} goal is generated and checked.

\section{Soundness Boundary and Trusted Computing Base}
\label{sec:tcb}

The metatheory depends on the following components:

\begin{enumerate}[leftmargin=1.4em,label=A\arabic*.]
  \item The C frontend and operational semantics agree on the supported
        fragment, including integer widths, structure layout, pointer offsets,
        and supported floating operations.
  \item Surface assertion elaboration is adequate for the accepted QCP
        annotation fragment.
  \item Canonicalization, exact encoding, and replay preserve typed blocks,
        aliases, roots, and concrete cell values.
  \item The finite statement, loop, and call rules are sound with respect to
        the supported C semantics and emit every required safety obligation.
  \item Whole-contract use of finite bindings is justified by a checked
        uniqueness, representative-adequacy, or ghost-completeness proof.
  \item Each normalization and strategy step is proof-producing or covered by
        a proved equivalence; final proofs target the original generated goal.
  \item SMT success appears in a proof module accepted by the final \rocq{}
        check; an SMT diagnostic alone is not evidence.
  \item The \rocq{} kernel and the whitelisted logical context are sound and
        consistent.
  \item The reference oracle correctly certifies positive and negative labels.
\end{enumerate}

The LLM, binding heuristics, mutation generator, test reducer, and proof search
policy are outside the trusted computing base.  They influence completeness
and cost, not the meaning of a checked verdict.

\begin{theorem}[End-to-end conditional soundness]
Under A1--A9, every \pass{} result establishes the expected candidate
pre/post proposition on the frozen concrete case, and every \fail{} result
establishes its negation.
\end{theorem}
\begin{proof}[Proof sketch]
By A1--A4, the execution artifacts denote the stored concrete states and
behavior, and the surface contract denotes the emitted internal assertion.
A5 connects a finite bind representation to the universal \texttt{With}
semantics.  A6--A8 and specialization preservation connect the checked closed
goal to that contract semantics.  The instance-verdict soundness theorem gives
the expected polarity.  A9 connects that polarity to the case's coverage or
precision label.
\end{proof}

This is a conditional system theorem, not yet a mechanized proof of the full
implementation.  QCP's checked proof aggregation directly supports A6--A8
for the VCs it emits; it does not automatically discharge A2--A5 for TeSpec's
new frontend/execution/evaluation path.  The artifact must identify which
assumptions are mechanized and which remain engineering trust.

\section{Limitations and Non-Claims}
\label{sec:limitations}

\paragraph{Finite tests are not universal verification.}
Passing every benchmark case is evidence relative to the suite.  It does not
imply extensional correctness outside the suite.  Mutation score measures
discrimination against a fault model, not proof of completeness.

\paragraph{Reference quality remains fundamental.}
An incorrect reference contract or mislabeled mutant corrupts the evaluation.
Both case polarities should be certified, and ambiguous cases should be
discarded rather than forced into a label.

\paragraph{Binding can remain ambiguous.}
Some heaps admit multiple valid abstract views.  Pre-state-only binding prevents
output leakage but does not imply semantic adequacy.  Without a uniqueness,
equivalence, or ghost-completeness proof, only the selected indexed instance
may receive a conclusive judgment; the whole-contract result is
\bindfail{}/\unknown.

\paragraph{C semantics are partial.}
Undefined behavior, implementation-defined operations, opaque external calls,
concurrency, and unsupported floating behavior must yield \unknown/\error{} or
use an explicit trusted model.  Concrete execution must not turn undefined
behavior into an arbitrary terminal state.

\paragraph{Inductive predicates may resist automation.}
An arbitrary \rocq{} predicate can be meaningful but impossible for QCP
strategies or SMT to unfold.  Such a case remains a residual proposition.
Supporting arbitrary syntax means preserving its semantics and returning
\unknown{} when necessary, not promising complete automation.

\paragraph{Proof search can bias resolution.}
Although conclusive verdicts are sound, candidates whose propositions are hard
to prove may receive more \unknown{} results.  Proof budgets, lower/upper
bounds, oracle proofs, and normalized-goal complexity must be reported.

\paragraph{Benchmark leakage and overfitting.}
Program-name-specific rewrites and fixed-corpus heuristics are prohibited.
Normalization is defined over typed assertion AST nodes.  Evaluation should
hold out projects, predicate families, and mutation families.

\section{Related Work}
\label{sec:related}

\paragraph{Specification-generation benchmarks.}
\textsc{Clever} checks candidate Lean specifications through equivalence to a
held-out reference and validates all artifacts with Lean's type
checker~\citep{thakur2025clever}.  FormalBench evaluates JML specification
consistency using deductive verification and approximates completeness with
mutation analysis~\citep{lecong2025formalbench}.  \verina{} provides modular
code, specification, and proof tasks in Lean and combines relation proofs with
concrete testing~\citep{ye2026verina}.  \coins{} specializes declarative
\rocq{} input/output propositions to concrete positive and negative behaviors,
directly motivating our instance-based design~\citep{yang2026coins}.
OSVBench evaluates state-machine specifications for operating-system
verification by comparing their verifier behavior over correct and buggy kernel
implementations~\citep{li2026osvbench}.  These works demonstrate complementary
ways to avoid human-only judgment.  \tespec{} focuses on the combination of
non-executable contracts, concrete mutable C heaps, separation ownership,
typed ghosts, and pre/post heap transitions.

\paragraph{Scaling test evidence.}
Property-based testing integrated with proof assistants, including
QuickChick and FuzzChick, provides executable generators and checkers with
strong metatheoretic foundations~\citep{paraskevopoulou2015quickchick,lampropoulos2019fuzzchick}.
\tespec{} reverses the trust direction: concrete behaviors are fixed, while the
possibly non-executable proposition is evaluated by proof.  VeriScale expands
and reduces positive/negative suites using adversarial
implementations~\citep{bai2026veriscale}; our heap-aware mutation protocol
adapts that principle to state transitions.

\paragraph{Heap logic and agentic synthesis.}
Separation-logic model checking establishes that rich fragments of symbolic
heaps with inductive predicates can be checked against concrete
states~\citep{brotherston2016model}.  \tespec{} combines this perspective with
program execution, full QCP assertions, and proof-assistant residuals.
Spec-Agent synthesizes first-order separation-logic specifications for large
C++ repositories using static analysis, heap traces, fuzzing, and agentic
refinement~\citep{suresh2026agentic}.  It shows that agents and heap
specifications can be combined at scale.  Our primary contribution is
different: the artifact under test is an existing/generated deductive
specification, while a deterministic proof boundary evaluates its acceptance
and rejection of fixed heap behaviors.

\section{Discussion: Why This Is Specification Evaluation}

It may appear that \tespec{} simply runs a C test.  Three distinctions make the
object of evaluation the specification.

First, the reference or mutant implementation only produces a labeled behavior.
The scored question is whether $P$ or $Q$ classifies that behavior as intended.
Second, the checked proposition is the original separation-logic assertion,
not an assertion compiled into ad hoc executable code.  Third, positive and
negative behavior generators can change while the candidate remains fixed;
the score measures the boundary drawn by the candidate relation.

The conceptual progression is:
\[
\underbrace{S(i,o)}_{\text{pure instantiated specification}}
\quad\longrightarrow\quad
\underbrace{\mathsf{Dom}_S(a,\sigma_0)\land
  \mathsf{Conf}_S(a,\sigma_0,r,\sigma_1)}_
 {\text{heap-transition contract instance}}.
\]
The additional components are not incidental engineering.  They create the
central research problems: fixing a candidate-independent heap, representing
heap isomorphism and aliasing, binding ghost state without output leakage,
executing calls and loops without turning the task into universal
verification, and retaining trustworthy verdicts when the logic is
non-executable.

\section{Conclusion}

\tespec{} generalizes instance-based formal-specification evaluation from pure
input/output values to C executions with mutable heaps.  A test is a frozen
state transition plus a certified adequate representation of any
pre-state-derived ghost parameters.  The target contract is elaborated and
specialized, not translated into an ad hoc executable property; finite loops
and visible callees are executed along the observed path; and both expected
and opposite logical polarities are checked.  The resulting three-valued
judgment keeps proof-search incompleteness visible and confines language
models to proposal roles.

The formalization also draws a necessary line between interactive contract
testing and comparative specification evaluation.  Materializing a heap from
the candidate precondition is convenient for the former, but a benchmark must
freeze and replay the same canonical heap for every candidate.  With that
discipline, certified positive and negative heap behaviors provide a scalable
way to measure the coverage and precision of expressive, non-executable
separation-logic specifications without claiming universal verification.

\bibliographystyle{plainnat}
\bibliography{references}

\end{document}
